%% mbeta_zero_LMP.tex
%% "ECI v7.4 sharp prediction m_1 = 0: falsifiability via neutrino
%%  mass ordering and future 0nbb searches"
%% Target: Letters in Mathematical Physics (LMP), 4-6 pp
%% Author: K. Remondiere (ECI v7.4, sub-agent M58/M77, 2026-05-06)
%%
%% Anti-hallucination policy:
%%   - NuFIT 6.0 values live-verified arXiv:2410.05380 (M58, 2026-05-06)
%%   - KamLAND-Zen 2025 PRL 135 live-verified arXiv:2406.11438
%%   - nEXO CDR live-verified arXiv:2106.16243
%%   - LEGEND-1000 live-verified arXiv:2107.11462
%%   - Tavartkiladze live-verified arXiv:2512.24804
%%   - Hallucination counter: 91 (no new fabrications introduced)
%%   - Mistral large-latest STRICT-BANNED

\documentclass[aps,prl,twocolumn,showpacs,superscriptaddress]{revtex4-2}

\usepackage{amsmath,amssymb,amsfonts}
\usepackage{hyperref}
\usepackage{xcolor}
\usepackage{booktabs}
\usepackage{bm}
\usepackage{url}

%% Shorthand macros
\newcommand{\Spf}{S'_4}
\newcommand{\tauS}{\tau_S = i}
\newcommand{\arXiv}[1]{\href{https://arxiv.org/abs/#1}{\texttt{arXiv:#1}}}
\newcommand{\mbb}{m_{\beta\beta}}
\newcommand{\Qi}{\mathbb{Q}(i)}

\begin{document}

\title{ECI v7.4 sharp prediction $m_1 = 0$: falsifiability via neutrino
       mass ordering and future $0\nu\beta\beta$ searches}

\author{K.~Remondi\`{e}re}
\affiliation{Independent researcher, Tarbes, France}
\email{kevin.remondiere@gmail.com}

\date{2026-05-06}

\begin{abstract}
The ECI v7.4 framework (NPP20 $S'_4$ modular symmetry with the
CSD$(1+\sqrt{6})$ Littlest Modular Seesaw at the fixed point
$\tau_S = i$) predicts a rank-2 Majorana mass matrix from a
two-right-handed-neutrino seesaw, forcing $m_1 = 0$ exactly in
normal ordering.  Using NuFIT~6.0 oscillation parameters
(arXiv:2410.05380), we compute the effective Majorana mass
$\mbb \in [1.50,\,3.72]$~meV over all Majorana phases.
This window lies \emph{below the sensitivity of all planned
neutrinoless double-beta-decay experiments} (KamLAND-Zen,
nEXO, LEGEND-1000), but is falsifiable via the neutrino mass
ordering: an inverted-ordering determination at $\geq 3\sigma$
by JUNO or DUNE-prec would refute the prediction.
We compare with the Tavartkiladze~\cite{Tav25} $S_3$ texture
(inverted ordering, $\mbb \gtrsim 34$~meV) and identify
mass-ordering determination as the decisive discriminant by
2030--2032.
\end{abstract}

\pacs{14.60.Pq, 23.40.Bw, 12.60.Jv}
\keywords{neutrinoless double beta decay, modular flavour symmetry,
  neutrino mass, seesaw, normal ordering}

\maketitle

%%--------------------------------------------------------------------
\section{Introduction}
\label{sec:intro}

The Extended Cosmological Index (ECI) research
programme~\cite{ECIv6053} anchors the lepton sector to the
LMFDB CM newform of label 4.5.b.a and selects the modular fixed
point $\tau_S = i$ via the CM-by-$\Qi$ structure.  At this
fixed point the $S'_4$ weight-2 triplet acquires the
alignment~\cite{DKLL19}
\begin{equation}
  Y_3^{(2)}(i) \;\propto\; (1,\;1+\sqrt{6},\;1-\sqrt{6}),
\end{equation}
which defines the CSD$(1+\sqrt{6})$ ansatz in the Littlest
Modular Seesaw (LMS)~\cite{LMS22}.  The two-right-handed-neutrino
(2-RHN) structure yields a rank-2 Dirac mass matrix, forcing
\begin{equation}
  m_1 = 0 \quad\text{(exact, Normal Ordering)},
  \label{eq:m1zero}
\end{equation}
with $m_2 = \sqrt{\Delta m^2_{21}}$,
$m_3 = \sqrt{\Delta m^2_{31}}$.

The prediction~\eqref{eq:m1zero} is sharp and structural: it
follows from the rank of the seesaw formula, not from a
fine-tuning.  The question then is whether $m_1 = 0$ is testable.
The most direct observable is the effective Majorana mass in
neutrinoless double-beta decay ($0\nu\beta\beta$),
$\mbb = |m_1 |U_{e1}|^2 + m_2 |U_{e2}|^2 e^{i\alpha}
        + m_3 |U_{e3}|^2 e^{i\beta}|$.
With $m_1 = 0$ this reduces to
\begin{equation}
  \mbb = \bigl| m_2\,|U_{e2}|^2\,e^{i\alpha_{21}}
              + m_3\,|U_{e3}|^2\,e^{i\alpha_{31}}\bigr|,
  \label{eq:mbb}
\end{equation}
where $\alpha_{21},\alpha_{31}$ are the two Majorana phases, free
within ECI v7.4 (the high-energy modular phase does not uniquely
pin the low-energy Majorana phases).

%%--------------------------------------------------------------------
\section{NuFIT 6.0 input and $\mbb$ window}
\label{sec:nufit}

We use the NuFIT~6.0 global fit~\cite{NuFIT60} (normal-ordering
best-fit values, without Super-Kamiokande atmospheric data):

\begin{table}[h]
\centering
\caption{NuFIT~6.0 Normal Ordering best-fit parameters used.}
\label{tab:nufit}
\begin{tabular}{lc}
\toprule
Parameter & Value \\
\midrule
$\sin^2\theta_{12}$ & $0.308$ \\
$\sin^2\theta_{13}$ & $0.02215$ \\
$\Delta m^2_{21}$   & $7.49\times10^{-5}\;\mathrm{eV}^2$ \\
$\Delta m^2_{31}$   & $+2.513\times10^{-3}\;\mathrm{eV}^2$ \\
\bottomrule
\end{tabular}
\end{table}

From Table~\ref{tab:nufit} one derives:
$m_2 = 8.654$~meV, $m_3 = 50.130$~meV,
$\Sigma m_\nu = 58.78$~meV.

The two amplitudes entering Eq.~\eqref{eq:mbb} are:
\begin{align}
  A_2 &\equiv m_2\,\sin^2\theta_{12}\,\cos^2\theta_{13}
       = 2.607\;\text{meV}, \label{eq:A2}\\
  A_3 &\equiv m_3\,\sin^2\theta_{13}
       = 1.110\;\text{meV}. \label{eq:A3}
\end{align}
The Majorana phase difference $\Delta\phi\equiv\alpha_{31}-\alpha_{21}$
scans $[0,2\pi)$, giving:
\begin{equation}
\boxed{
  \mbb \;\in\; [1.50,\;3.72]\;\text{meV}
  \quad\text{(NuFIT 6.0, all Majorana phases)}
}
\label{eq:mbb_range}
\end{equation}
with maximum at $\Delta\phi = 0$ ($A_2 + A_3 = 3.72$~meV) and
minimum at $\Delta\phi = \pi$ ($|A_2 - A_3| = 1.50$~meV).

%%--------------------------------------------------------------------
\section{Experimental landscape}
\label{sec:landscape}

Table~\ref{tab:expt} summarises the relevant $0\nu\beta\beta$
experimental program and its sensitivity relative to the ECI
prediction~\eqref{eq:mbb_range}.

\begin{table}[h]
\centering
\caption{$0\nu\beta\beta$ experimental landscape vs.\ ECI v7.4
prediction $\mbb\in[1.50,3.72]$~meV.
``NO'' = not observable; ``M'' = marginal (optimistic NME).}
\label{tab:expt}
\begin{tabular}{lccc}
\toprule
Experiment & Isotope & $\mbb$ sensitivity & ECI \\
\midrule
KamLAND-Zen 2025~\cite{KLZ25}   & $^{136}$Xe & $28$--$122$~meV & NO \\
KamLAND-Zen 2027 proj.           & $^{136}$Xe & $\sim12$--$20$~meV & NO \\
nEXO 2030+~\cite{nEXO21}        & $^{136}$Xe & $5$--$20$~meV   & M \\
LEGEND-1000 2030+~\cite{LEG21}  & $^{76}$Ge  & $9$--$21$~meV   & NO \\
Post-2035 ($<2$~meV)            & TBD        & $<2$~meV         & decisive \\
\bottomrule
\end{tabular}
\end{table}

The ECI window [1.50, 3.72]~meV is a factor 3--13 below the
lower edge of nEXO optimistic sensitivity.  To reach $\mbb\sim3.72$~meV
requires $T_{1/2}\gtrsim 2\times10^{28}$~yr; to reach
$\mbb\sim1.50$~meV requires $T_{1/2}\gtrsim10^{29}$~yr (NME-dependent).

A positive signal at $\mbb > 5$~meV in any current-generation
experiment would imply $m_1\neq 0$, \emph{falsifying}
the rank-2 structure.  A signal at $\mbb > 15$~meV would further
require inverted ordering, doubly falsifying ECI v7.4.

%%--------------------------------------------------------------------
\section{Indirect falsifiability: mass ordering}
\label{sec:falsify}

Since direct $0\nu\beta\beta$ detection is out of reach for
near-term experiments, the primary falsifier is the \emph{neutrino
mass ordering}:

\begin{enumerate}
\item[\textbf{F8d}] \textbf{Inverted ordering at $\geq3\sigma$
  (JUNO 2030, DUNE-prec 2032+).}
  ECI v7.4 demands normal ordering from the 2-RHN seesaw
  ($m_1=0$ is inconsistent with inverted ordering, which requires
  $m_3 = 0$ or quasi-degenerate spectrum).  An IO determination at
  $\geq 3\sigma$ by JUNO~\cite{JUNO22} or DUNE~\cite{DUNE20}
  \emph{falsifies} the lepton sector of ECI v7.4.

\item[\textbf{F8a}] \textbf{$\delta_{CP}\neq -87^\circ$ at
  $\geq3\sigma$ (DUNE 2032+).}
  The LMS at $\tau_S=i$ predicts $\delta_{CP}\approx-87^\circ$.
  DUNE precision ($\pm15^\circ$) provides an independent
  lepton-sector falsifier.

\item[\textbf{F8b/c}] \textbf{$\mbb > 5$~meV detected.}
  Falsifies rank-2 ($m_1\neq0$).
\end{enumerate}

Null results at KamLAND-Zen / nEXO / LEGEND-1000 are
\emph{consistent} with, but not confirmatory of, the ECI
prediction.  Confirmation requires either (i)~a future experiment
reaching $\mbb\sim1$--$3$~meV, or (ii)~a joint global fit
confirming normal ordering with $m_1 \lesssim 1$~meV at high
significance.

%%--------------------------------------------------------------------
\section{Comparison with Tavartkiladze (2025)}
\label{sec:compare}

The Tavartkiladze 2025 model~\cite{Tav25} employs a
$\Gamma_2\simeq S_3$ modular symmetry (level-2, non-holomorphic)
with a distinct texture that favours inverted ordering and
quasi-degenerate masses, predicting $\mbb\gtrsim34$~meV.

Table~\ref{tab:compare} summarises the key differences.

\begin{table}[h]
\centering
\caption{ECI v7.4 vs.\ Tavartkiladze 2025~\cite{Tav25}.}
\label{tab:compare}
\begin{tabular}{lcc}
\toprule
Property & ECI v7.4 & Tavartkiladze \\
\midrule
Modular group  & $S'_4$ (level-4) & $S_3$ (level-2) \\
Fixed point    & $\tau_S = i$     & distinct \\
Ordering       & Normal           & Inverted \\
$m_1$          & $0$ (exact)      & non-zero \\
$\mbb$         & $[1.50,3.72]$~meV & $\gtrsim34$~meV \\
\bottomrule
\end{tabular}
\end{table}

The two predictions differ by more than an order of magnitude.
If JUNO/DUNE establish normal ordering at $\geq3\sigma$: ECI v7.4
is corroborated, Tavartkiladze constrained.
If inverted ordering is established: ECI v7.4 is falsified,
Tavartkiladze corroborated.  The discriminant is therefore the
mass-ordering determination, well within reach by 2030--2032.

%%--------------------------------------------------------------------
\section{Discussion and conclusion}
\label{sec:concl}

The ECI v7.4 combination (NPP20 $S'_4$, LMS at $\tau_S=i$,
CSD$(1+\sqrt{6})$) yields the sharp prediction $m_1=0$ from
rank-2 seesaw structure, giving $\mbb\in[1.50,3.72]$~meV
with NuFIT~6.0 parameters.  This window is:
\begin{itemize}
\item below all near-term $0\nu\beta\beta$ experimental sensitivity;
\item consistent with current null results;
\item falsifiable by mass-ordering experiments (JUNO/DUNE, 2030--2032);
\item qualitatively distinct from the Tavartkiladze IO prediction.
\end{itemize}

We emphasise that the title ``falsifiability via neutrino mass
ordering'' is precise: the ordering determination is the
experimentally accessible falsifier, not direct $0\nu\beta\beta$
detection.  A post-2035 experiment reaching $T_{1/2}\sim10^{29}$~yr
would provide the definitive $\mbb$ test.

%%--------------------------------------------------------------------
\section*{Acknowledgements}

The author thanks the LMFDB collaboration for the modular form database
that anchors the analysis, the NuFIT collaboration for the publicly
available global oscillation fits, and the KamLAND-Zen, JUNO, and
DUNE collaborations for sensitivity projections.

%%--------------------------------------------------------------------
\begin{thebibliography}{99}

\bibitem{ECIv6053}
K.~Remondi\`ere,
\emph{ECI/crossed-cosmos repository, version 6.0.53},
Zenodo (2026-05-05);
DOI \href{https://doi.org/10.5281/zenodo.20036808}{10.5281/zenodo.20036808}.

\bibitem{DKLL19}
G.-J.~Ding, S.~F.~King, X.-G.~Liu, J.-N.~Lu,
\emph{Modular $S_4$ and $A_4$ symmetries and their fixed points:
new predictive examples of lepton mixing},
JHEP \textbf{12} (2019) 030;
\arXiv{1910.03460}.
%% [VERIFIED 2026-05-05 via arXiv API.]

\bibitem{LMS22}
I.~de~Medeiros~Varzielas, S.~F.~King, M.~Levy,
\emph{Littlest Modular Seesaw},
\arXiv{2211.00654} (2022).
%% [VERIFIED 2026-05-05 via arXiv API.]

\bibitem{NPP20}
P.~P.~Novichkov, J.~T.~Penedo, S.~T.~Petcov,
\emph{Double Cover of Modular $S_4$ for Flavour Model Building},
Nucl.~Phys.~B \textbf{963} (2021) 115301;
\arXiv{2006.03058}.
%% [VERIFIED 2026-05-05 via arXiv API.]

\bibitem{NuFIT60}
I.~Esteban, M.~C.~Gonz\'alez-Garc\'ia, M.~Maltoni,
T.~Schwetz, A.~Zhou,
\emph{Nu-fit 6.0: new results for neutrino oscillation parameters
from global analysis of current data},
\arXiv{2410.05380} (2024).
%% [LIVE-VERIFIED 2026-05-06 by M58 via WebFetch.]

\bibitem{KLZ25}
S.~Abe et al.\ (KamLAND-Zen Collaboration),
\emph{Search for the Majorana Nature of Neutrinos in the Inverted
Mass Ordering Region with KamLAND-Zen},
Phys.~Rev.~Lett.\ \textbf{135} (2025) 262501;
\arXiv{2406.11438}.
%% [LIVE-VERIFIED 2026-05-06 by M58 via WebFetch.]

\bibitem{nEXO21}
G.~Adhikari et al.\ (nEXO Collaboration),
\emph{nEXO: neutrinoless double beta decay search beyond
$10^{28}$ year half-life sensitivity},
J.~Phys.~G \textbf{49} (2022) 015104;
\arXiv{2106.16243}.
%% [LIVE-VERIFIED 2026-05-06 by M58 via WebFetch.]

\bibitem{LEG21}
N.~Abgrall et al.\ (LEGEND Collaboration),
\emph{LEGEND-1000 Preconceptual Design Report},
\arXiv{2107.11462} (2021).
%% [LIVE-VERIFIED 2026-05-06 by M58 via WebFetch.]

\bibitem{JUNO22}
A.~Abusleme et al.\ (JUNO Collaboration),
\emph{JUNO Physics and Detector},
Prog.\ Part.\ Nucl.\ Phys.\ \textbf{123} (2022) 103927;
\arXiv{2204.13249}.
%% [VERIFIED 2026-05-05 via arXiv API.]

\bibitem{DUNE20}
B.~Abi et al.\ (DUNE Collaboration),
\emph{Long-baseline neutrino oscillation physics potential of the
DUNE experiment},
Eur.~Phys.~J.~C \textbf{80} (2020) 978;
\arXiv{2006.16043}.
%% [VERIFIED 2026-05-05 via arXiv API.]

\bibitem{Tav25}
G.~Tavartkiladze,
\emph{Minimal Modular Flavor Symmetry and Lepton Textures Near Fixed Points},
\arXiv{2512.24804} (2025).
%% [LIVE-VERIFIED 2026-05-06 by M58 via WebFetch.]

\end{thebibliography}

\end{document}
